\SourceLeaf{19}
\begin{center}
{\large\bfseries SECTION TWO}\par\bigskip
{\small\bfseries ON}\par\bigskip
{\large\bfseries CONGRUENCES OF THE FIRST DEGREE.}
\end{center}
\begin{center}\rule{0.09\linewidth}{0.3pt}\end{center}
\begin{center}\emph{Preliminary theorems on prime numbers, factors, etc.}\end{center}
\Article{13}
\textsc{Theorem}. \emph{The product of two positive numbers smaller than a given prime number is not divisible by that prime.}

Let $p$ be prime, and $a$ positive and $<p$: then there will be no positive number $b$ smaller than $p$ such that $ab\equiv 0\ (\mathrm{mod.}\ p)$.

\emph{Proof.} If anyone denies this, let us suppose that there are numbers $b, c, d$ etc., all $<p$, such that $ab\equiv 0$, $ac\equiv 0$, $ad\equiv 0$ etc. (mod. $p$). Let $b$ be the smallest of all, so that all numbers smaller than $b$ lack this property. Clearly $b>1$: for if $b=1$, we would have $ab=a<p$ (\emph{hyp.}), so it would not be divisible by $p$. So $p$, being prime, cannot be divided by $b$, but will fall between the two nearest multiples of $b$, $mb$ and $(m+1)b$. Let $p-mb=b'$; then $b'$ will be a positive number and $<b$. Now, since we have supposed that $ab\equiv 0\ (\mathrm{mod.}\ p)$, we also have $mab\equiv 0$ (art. 7), and so, subtracting this from $ap\equiv 0$, we obtain $a(p-mb)=ab'\equiv 0$; that is, $b'$ must be included among the numbers $b, c, d$ etc., although it is smaller than the least of them, $b$. \emph{Q.~E.~A.}

\Article{14}
\emph{If neither $a$ nor $b$ can be divided by the prime number $p$: their product $ab$ likewise cannot be divided by $p$.}

\SourceLeaf{20}
Let the least positive residues of the numbers $a, b$, with respect to the modulus $p$, be $\alpha, \beta$, neither of which will be $0$ (\emph{hyp.}) Now if $ab\equiv 0$ (mod. $p$), then, because $ab\equiv\alpha\beta$, we would also have $\alpha\beta\equiv 0$, which is incompatible with the preceding theorem.

A proof of this theorem was already given by Euclid, \emph{El.} VII. 32. Nevertheless, we did not wish to omit it, both because several more recent writers have either passed off vague reasoning as a proof or omitted the theorem altogether, and because this simpler case makes it easier to understand the method used here, which we will use later to resolve much deeper matters.

\Article{15}
\emph{If none of the numbers $a, b, c, d$ etc. can be divided by the prime number $p$, their product $abcd$ etc. likewise cannot be divided by $p$.}

By the preceding article, $ab$ cannot be divided by $p$; therefore neither can $abc$; hence neither can $abcd$ etc.

\Article{16}
\textsc{Theorem}. \emph{Any composite number can be decomposed into prime factors in only one way.}

\emph{Proof.} That every composite number can be decomposed into prime factors is known from the elements, but it is generally, and improperly, assumed without saying so that this cannot be done in several different ways. Let us suppose that a composite number $A$, which is $=a^\alpha b^\beta c^\gamma$ etc., where $a, b, c$ etc. denote distinct prime numbers, can be decomposed into prime factors in yet another way. First, it is evident that no primes other than $a, b, c$ etc. can enter this second system of factors, since any other prime cannot divide the number $A$ composed of these. Likewise, none of the primes $a, b, c$ etc. can be absent from this second system of factors, since otherwise it would not divide $A$ (preceding art.). Thus these two decompositions into factors can differ only in that some prime occurs more often in one than in the other. Let $p$ be such a prime, occurring $m$ times in one decomposition, but $n$ times in the other, and let $m>n$: now remove the factor $p$ from each system $n$ times, so that in one it still remains $m-n$ times, while from the other it has entirely disappeared. That is, there are two decompositions of the number $\frac{A}{p^n}$ into factors, one of which is wholly free of the factor $p$, while the other contains it $m-n$ times, contrary to what we have just proved.

\SourceLeaf{21}
\Article{17}
If, therefore, a composite number $A$ is the product of $B, C, D$ etc., it is clear that among the prime factors of the numbers $B, C, D$ etc. there can be none except those that are also among the factors of the number $A$, and that each of these factors must occur in $B, C, D$ etc., taken together, as often as in $A$. From this we obtain a criterion for whether the number $B$ divides another number $A$ or not. It will do so if $B$ contains neither other prime factors nor any one of them more often than $A$; if either of these conditions fails, $B$ will not divide $A$.

It can easily be derived from this, with the aid of counting combinations, that if $A=a^\alpha b^\beta c^\gamma$ etc., where, as above, $a, b, c$ etc. denote distinct prime numbers: $A$ has
\[
(\alpha+1)(\beta+1)(\gamma+1)\ \mathrm{etc.}
\]
distinct divisors, including $1$ and $A$.

\Article{18}
If, then, $A=a^\alpha b^\beta c^\gamma$ etc., $K=k^\kappa l^\lambda m^{\mu}$ etc., and the primes $a, b, c$ etc., $k, l, m$ etc. are all distinct, it is clear that $A$ and $K$ have no common divisor except $1$, or are coprime.

Given several numbers $A, B, C$ etc., their \emph{greatest common divisor} is determined as follows. Decompose them all into their prime factors, and select from these those that are common to all the numbers $A, B, C$ etc. (if there are no such factors, no divisor will be common to all). Then record how often each of these prime factors is contained in each of $A, B, C$ etc., or \emph{what exponent} each has in the individual numbers $A, B, C$ etc. Finally, assign to each prime factor the least of all the exponents it has in $A, B, C$ etc., and form their product, which will be the common measure sought.

But when the \emph{least common multiple} of the numbers $A, B, C$ etc. is desired, proceed as follows. Collect all the primes that divide any one of the numbers $A, B, C$ etc., assign to each the greatest of all the exponents it has in the numbers $A, B, C$ etc., and thus form the product of them all, which will be the multiple sought.

\emph{Example.} Let $A=504=2^3 3^2 7$; $B=2880=2^6 3^2 5$; $C=864=2^5 3^3$. To find the greatest common divisor we have the prime factors $2, 3$, to which the exponents $3, 2$ are to be assigned; hence the result is $=2^3 3^2=72$; but the least common multiple will be $2^6 3^3 5.7=60480$.

\SourceLeaf{22}
We omit the proofs because they are easy. Moreover, how these problems are to be solved when the resolution of the numbers $A, B, C$ etc. into factors is not given is known from the elements.

\Article{19}
\emph{If the numbers $a, b, c$ etc. are coprime to another number $k$, their product $abc$ etc. is also coprime to $k$.}

For since none of the numbers $a, b, c$ etc. has a prime factor in common with $k$, and the product $abc$ etc. cannot have prime factors other than those that are factors of one of the numbers $a, b, c$ etc., the product $abc$ etc. will likewise have no prime factor in common with $k$. So, by the preceding article, $k$ is coprime to $abc$ etc.

\emph{If the numbers $a, b, c$ etc. are pairwise coprime, and each divides another number $k$: their product will also divide the number $k$.}

This follows just as easily from arts. 17, 18. For let any prime divisor of the product $abc$ etc. be $p$, occurring in it $\pi$ times, and it is clear that one of the numbers $a, b, c$ etc. must contain this same divisor $\pi$ times. Therefore $k$ also, which this number divides, contains $\pi$ copies of the divisor $p$. The same holds for the remaining divisors of the product $abc$ etc.

So \emph{if two numbers $m, n$ are congruent with respect to several moduli $a, b, c$ etc. that are pairwise coprime, they will also be congruent with respect to their product.} For since $m-n$ is divisible by each of $a, b, c$ etc., it can also be divided by their product.

Finally, \emph{if $a$ is coprime to $b$ and $ak$ is divisible by $b$, then $ak$ will also be divisible by $b$.} For since $ak$ is divisible both by $a$ and by $b$, it can also be divided by $ab$; that is, $\frac{ak}{ab}=\frac{k}{b}$ will be an integer.

\Article{20}
\emph{When $A=a^\alpha b^\beta c^\gamma$ etc., where $a, b, c$ etc. denote distinct prime numbers, is some power, say $=k^n$: all the exponents $\alpha, \beta, \gamma$ etc. will be divisible by $n$.}

For the number $k$ contains no prime factors other than $a, b, c$ etc. Let it contain the factor $a$ exactly $\alpha'$ times; then $k^n$, or $A$, will contain this factor $n\alpha'$ times; so $n\alpha'=\alpha$, and $\frac{\alpha}{n}$ is an integer. Similarly, $\frac{\beta}{n}$ etc. are shown to be integers.
